\section{Stopping sets survive positive raw cancellation}
\label{sec:stopping}

\begin{definition}[Positive raw pivot]
Let $N=(n_{ih})$ be a nonnegative integral matrix. A positive raw pivot at
$(r,g)$ is allowed when $n_{rg}=1$. It deletes row $r$ and column $g$ and
replaces the other entries by \eqref{eq:positive-pivot}. A full raw
cancellation is a finite sequence of such pivots that removes every row.
\end{definition}

The adjective positive refers to the plus sign in the update. In particular,
this operation is not Gaussian elimination over a field. It records a
geometric disk-slide model before any subsequent reduction of intersections.

\begin{definition}[Stopping set]
A nonempty set $R$ of rows of $N$ is a stopping set if, for each column $h$,
\begin{equation}\label{eq:stopping}
 S_h(R):=\sum_{i\in R}n_{ih}\in\{0,2,3,\ldots\}.
\end{equation}
Thus no column has exactly one crossing into $R$. The empty set is admitted
when discussing unions, but it is never an obstruction certificate.
\end{definition}

A coefficient two in a single row counts as two crossings. Replacing the
matrix by its Boolean support would lose precisely this distinction.

\begin{theorem}[Persistence under every positive pivot]
\label{thm:persistence}
Suppose $R$ is a stopping set of $N$ and a positive raw pivot is performed
at $(r,g)$. If $r\notin R$, then $R$ is a stopping set of the new matrix.
If $r\in R$, then $R\setminus\{r\}$ is nonempty and is a stopping set of
the new matrix. Consequently a matrix with a nonempty stopping set admits
no full raw cancellation.
\end{theorem}

\begin{proof}
First let $r\notin R$. For every surviving column $h$,
\[
 S'_h(R)=S_h(R)+S_g(R)n_{rh}.
\]
Each summand is either zero or at least two, because $S_h(R)$ and $S_g(R)$
have that property. Their sum therefore cannot equal one.

Now let $r\in R$. Since $n_{rg}=1$, the stopping condition at $g$ forces
$S_g(R)\ge2$. Put
\[
 k=\sum_{i\in R\setminus\{r\}}n_{ig}=S_g(R)-1\ge1.
\]
In particular $R\setminus\{r\}$ is nonempty. For any surviving $h$,
\begin{align}
 S'_h(R\setminus\{r\})
 &=S_h(R)-n_{rh}+k n_{rh}\notag\\
 &=S_h(R)+(k-1)n_{rh}.\label{eq:stopping-persistence}
\end{align}
If $S_h(R)=0$, nonnegativity implies $n_{rh}=0$ and the expression is zero.
If $S_h(R)\ge2$, it remains at least two. This proves persistence in both
cases. Inductively, a nonempty stopping set survives every prefix of a
positive pivot sequence, so such a sequence cannot remove every row.
\end{proof}

The short identity \eqref{eq:stopping-persistence} is useful because it
rules out \emph{all} positive pivot orders directly. It does not presume
a matching, a chosen survivor set or a deterministic producer.

\begin{definition}[Column-leaf deletion]
A column $g$ is a unit leaf in the current matrix if
$\sum_i n_{ig}=1$. There is then a unique incident row $r$, its entry is
one, and its positive pivot has no fill: all surviving $n_{ig}$ vanish.
A column-leaf deletion removes this row and column.
\end{definition}

\begin{theorem}[Complete rectangular raw cancellation]
\label{thm:complete}
For every finite nonnegative integral matrix $N$, the following are equivalent:
\begin{enumerate}[label=(\roman*)]
\item $N$ admits a full positive raw cancellation;
\item $N$ has no nonempty stopping set;
\item every maximal sequence of unit column-leaf deletions removes all rows;
\item some sequence of unit column-leaf deletions removes all rows.
\end{enumerate}
In particular, deciding full raw cancellation requires no branching over
pivot orders or over the columns that will survive.
\end{theorem}

\begin{proof}
Theorem~\ref{thm:persistence} gives (i)$\Rightarrow$(ii).
Run any maximal column-leaf sequence. If rows remain, call their set $R$.
Every undeleted column has total zero or at least two on $R$, because the
sequence is maximal. Each deleted column had only its designated row live
when it was removed; hence its total on the final $R$ is zero. No fill has
ever occurred. Therefore $R$ is a stopping set of the original matrix.
This proves (ii)$\Rightarrow$(iii). Implication (iii)$\Rightarrow$(iv) is
immediate, and every leaf deletion is a valid raw pivot, giving
(iv)$\Rightarrow$(i).
\end{proof}

\begin{corollary}[A canonical obstruction]
\label{cor:canonical}
There is a unique largest stopping set, with the empty set allowed.
Every maximal column-leaf sequence leaves exactly this row set, denoted
$\Stop(N)$. Thus the residual row set is independent of tie choices,
although pivot columns and successful schedules need not be unique.
\end{corollary}

\begin{proof}
The union of stopping sets is a stopping set: if a column meets the union,
it meets one constituent set and already has total at least two there.
Because there are finitely many rows, their union gives the largest set.
No column leaf can delete a row in an extant stopping set, since that set
would have column sum one. Hence every original stopping set lies in the
final residual. The proof of Theorem~\ref{thm:complete} shows that the
residual is itself a stopping set. The two containments prove the claim.
\end{proof}

This canonical peeling statement is the classical stopping-set phenomenon
identified in the prior-art discussion. The additional persistence theorem
explains why the same obstruction remains valid when the allowed operations
are enlarged from leaf deletions to all positive raw pivots.

\subsection{Exact saturation and sparse representation}

\begin{lemma}[Saturation at two is exact for the model]
\label{lem:saturation}
The map $\capTwo(a)=\min(2,a)$ preserves addition and multiplication when
the target operations are capped at two. It preserves the singleton
predicate and the stopping-set predicate. It therefore preserves the
existence of a full positive raw cancellation.
\end{lemma}

\begin{proof}
For $a,b\in\NN$, both
$\min(2,a+b)=\min(2,\min(2,a)+\min(2,b))$ and
$\min(2,ab)=\min(2,\min(2,a)\min(2,b))$ follow by splitting into whether
the result is zero, one or at least two. The same identities commute
with \eqref{eq:positive-pivot}. A sum is one exactly when one summand is
one and all the others are zero, which is also unchanged by capping.
\end{proof}

Actual crossing multiplicities may consequently be arbitrarily large in
binary without changing the finite peeling state. Values zero, one and
at least two are all that the decision needs. This statement does not
authorize replacing a geometric count by a signed exponent sum.

\subsection{A matching interpretation, without a matching algorithm}

\begin{corollary}[Unit uniquely restricted row matching]
\label{cor:matching}
A full raw cancellation exists if and only if there is a set of $m$
columns whose $m\times m$ submatrix has permanent one. Equivalently,
the weighted bipartite incidence multigraph has a row-saturating matching
whose induced square subgraph has a unique perfect matching, all of whose
chosen edge multiplicities are one.
\end{corollary}

\begin{proof}
Suppose peeling removes all rows. Select the $m$ pivot columns and order
both rows and those columns by deletion time. A pivot column has its
diagonal entry one and no nonzero in any later row. Thus the selected
square matrix is triangular with unit diagonal. Its only nonzero
permanent term is the diagonal term.

Conversely, a nonnegative integral square matrix of permanent one has a
unique perfect matching with unit chosen entries. Orient an off-matching
incidence from its row's matched pair to the incident column's matched
pair. A directed cycle would give an alternating cycle and a second
perfect matching. The directed graph is therefore acyclic and has a
source. Its matched column has no incidence from other rows. Removing
that pair and repeating gives column-leaf deletion of the selected minor.
These are also column leaves in the full matrix, because the selected
minor includes every row. Other columns do not affect a leaf's column sum.
\end{proof}

There is no need to compute a permanent or to enumerate the selected
columns. The direct obstruction theorem supplies the complete algorithm.
The distinction between \emph{saturating all rows} and maximizing the
size of a partially constrained matching should be retained when comparing
this corollary with matching literature.

\subsection{Boundary cases and limits of the conclusion}

\begin{example}[A zero row]
A single all-zero row is a stopping set. A raw singleton cancellation
cannot remove its two-handle, even if that handle has some other geometric
disposal, such as a later two/three-handle cancellation. Three-handles are
outside the present model.
\end{example}

\begin{example}[Peeling is not maximum partial cancellation]
For
$N=\left(\begin{smallmatrix}1&1\\1&1\end{smallmatrix}\right)$, every
column has total two and no column leaf exists. A positive pivot can
nevertheless remove one row and leave the matrix $(2)$. Thus
$|\Stop(N)|$ is not the minimum number of rows attainable by an arbitrary
partial raw sequence. The theorem decides complete cancellation and gives
a canonical source obstruction; it does not optimize partial progress.
\end{example}

\begin{example}[A unimodular raw obstruction]\label{ex:unimodular}
The matrix
\[
 U=\begin{pmatrix}1&1\\1&2\end{pmatrix}
\]
has determinant one, but its entire row set is a stopping set. Every legal
first raw pivot leaves the single entry three. Signed linear algebra
therefore cannot replace the unsigned criterion. The group presentation
$\langle x,y\mid xy,xy^2\rangle$ is trivial, since the two equations give
$y=1$ and then $x=1$. Its raw occurrence matrix is $U$.

There is also a geometric ball realization. Begin with the standard
genus-two handlebody and its two cancelling two-handles, attached along
the two primitive generator curves. Slide the first attaching curve over
the second to obtain $xy$, then slide the second over the first to obtain
the cyclic word $xy^2$, choosing the bands with no additional cocore crossings. These are
ordinary geometric two-handle slides, so the resulting manifold is still
a ball. The indicated raw incidence is $U$. This supplies a concrete
reason that a raw obstruction must not be reported as a nonball or
knottedness certificate.
\end{example}
